\documentclass[10pt,a4paper]{article}
\usepackage[utf8]{inputenc}
\usepackage[T5]{fontenc}
\usepackage[vietnamese]{babel}
\usepackage[margin=1.3cm,top=0.78cm,bottom=1.02cm,includehead,includefoot,headheight=20pt,headsep=6pt,footskip=16pt]{geometry}
\usepackage{amsmath,amssymb,array,tabularx,enumitem,fancyhdr,tikz,xcolor,needspace,multicol,colortbl,booktabs,microtype,tcolorbox}
\usepackage{tgheros}
\usetikzlibrary{calc,angles,quotes,arrows.meta,babel}
\tcbuselibrary{skins}

% UI/UX Pro Max Design Tokens: Minimalism & Swiss Style + Education Teal Palette
\definecolor{ink}{HTML}{134E4A}
\definecolor{brand}{HTML}{0D9488}
\definecolor{brandDark}{HTML}{0F766E}
\definecolor{accent}{HTML}{0F766E}
\definecolor{brandTint}{HTML}{F0FDFA}
\definecolor{brandLine}{HTML}{99F6E4}
\definecolor{slateText}{HTML}{334155}
\definecolor{mutedSlate}{HTML}{64748B}
\definecolor{ruleGray}{HTML}{CBD5E1}
\definecolor{amberAccent}{HTML}{D97706}
\definecolor{amberTint}{HTML}{FFFBEB}

\setlength{\parindent}{0pt}
\setlength{\parskip}{0pt}
\setlength{\columnseprule}{0.5pt}
\def\columnseprulecolor{\color{ruleGray}}

\newcommand{\headerQuote}{%
  \ifcase\value{page}\relax
    \or Math gives us hope that every problem has a solution% Trang 1
    \or Math + You = Champion% Trang 2
    \or Where Frustration Becomes Inspiration% Trang 3
    \or Turning Confusion into Clarity% Trang 4
    \or Don't be an X, be a Y of mathematics% Trang 5
    \or Don't get mad. Learn to add% Trang 6
    \or The only way to learn mathematics is to do mathematics% Trang 7
    \or Small steps every day lead to big results% Trang 8
    \or Every expert was once a beginner% Trang 9
    \or Success is the sum of small efforts repeated daily% Trang 10
    \or Mathematics is the music of reason% Trang 11
    \or Believe you can and you're halfway there% Trang 12
    \or Math gives us hope that every problem has a solution% Trang 13
    \or Math + You = Champion% Trang 14
    \or Where Frustration Becomes Inspiration% Trang 15
    \or Turning Confusion into Clarity% Trang 16
    \or Don't be an X, be a Y of mathematics% Trang 17
    \or Don't get mad. Learn to add% Trang 18
    \or The only way to learn mathematics is to do mathematics% Trang 19
    \or Small steps every day lead to big results% Trang 20
    \or Every expert was once a beginner% Trang 21
    \or Success is the sum of small efforts repeated daily% Trang 22
    \or Mathematics is the music of reason% Trang 23
    \or Believe you can and you're halfway there% Trang 24
    \or Math gives us hope that every problem has a solution% Trang 25
    \or Math + You = Champion% Trang 26
    \or Where Frustration Becomes Inspiration% Trang 27
    \or Turning Confusion into Clarity% Trang 28
    \or Don't be an X, be a Y of mathematics% Trang 29
    \or Don't get mad. Learn to add% Trang 30
    \else Math gives us hope that every problem has a solution%
  \fi
}

\pagestyle{fancy}
\fancyhf{}
\lhead{\footnotesize\textsf{\textit{\color{mutedSlate}\headerQuote}}}
\rhead{\footnotesize\textsf{\textbf{\color{ink}Người biên soạn:} \color{slateText}Dương Gia Bảo}}
\lfoot{\footnotesize\textsf{\color{mutedSlate}Toán 10 KNTT \textbullet\ Chương IV. Vectơ}}
\rfoot{\tikz[baseline=(p.base)]\node[fill=brandTint,draw=brandLine,rounded corners=2.5pt,inner xsep=5pt,inner ysep=1.8pt,font=\footnotesize\sffamily\bfseries,text=ink] (p) {Trang \thepage/\,3};}
\renewcommand{\headrulewidth}{0.5pt}
\renewcommand{\headrule}{\hbox to\headwidth{\color{ruleGray}\leaders\hrule height \headrulewidth\hfill}}
\renewcommand{\footrulewidth}{0pt}

\fancypagestyle{dapan}{
  \fancyhf{}
  \lhead{\footnotesize\textsf{\textit{\color{mutedSlate}\headerQuote}}}
  \rhead{\footnotesize\textsf{\textbf{\color{ink}Người biên soạn:} \color{slateText}Dương Gia Bảo}}
  \lfoot{\footnotesize\textsf{\textit{\color{mutedSlate}Đáp án Đề ôn tập Chương IV. Vectơ (Đề số 6) — Toán 10 KNTT}}}
  \rfoot{\tikz[baseline=(p.base)]\node[fill=brandTint,draw=brandLine,rounded corners=2.5pt,inner xsep=5pt,inner ysep=1.8pt,font=\footnotesize\sffamily\bfseries,text=ink] (p) {Trang \thepage/\,2};}
  \renewcommand{\headrulewidth}{0.5pt}
  \renewcommand{\headrule}{\hbox to\headwidth{\color{ruleGray}\leaders\hrule height \headrulewidth\hfill}}
  \renewcommand{\footrulewidth}{0pt}
}

\newcommand{\parttitle}[1]{%
  \par\vspace{4.5pt}\needspace{3\baselineskip}%
  \begin{tcolorbox}[enhanced,frame hidden,interior style={left color=brandTint,right color=white},borderline west={3pt}{0pt}{brand},sharp corners,left=6pt,right=4pt,top=3pt,bottom=3pt,boxsep=0pt,before skip=0pt,after skip=2pt]
    {\sffamily\normalsize\bfseries\color{ink}#1}
  \end{tcolorbox}\nopagebreak%
}
\newcommand{\four}[4]{\par\vspace{1.2pt}\noindent\begin{tabularx}{\linewidth}{@{}*{4}{>{\raggedright\arraybackslash}X}@{}}\textsf{\textbf{\color{ink}A.}} #1&\textsf{\textbf{\color{ink}B.}} #2&\textsf{\textbf{\color{ink}C.}} #3&\textsf{\textbf{\color{ink}D.}} #4\end{tabularx}}
\newcommand{\two}[4]{\par\vspace{1.2pt}\noindent\begin{tabularx}{\linewidth}{@{}*{2}{>{\raggedright\arraybackslash}X}@{}}\textsf{\textbf{\color{ink}A.}} #1&\textsf{\textbf{\color{ink}B.}} #2\\[3.5pt]\textsf{\textbf{\color{ink}C.}} #3&\textsf{\textbf{\color{ink}D.}} #4\end{tabularx}}
\newcommand{\one}[4]{\par\vspace{1.2pt}\noindent\begin{tabularx}{\linewidth}{@{}>{\raggedright\arraybackslash}X@{}}\textsf{\textbf{\color{ink}A.}} #1\\[3.5pt]\textsf{\textbf{\color{ink}B.}} #2\\[3.5pt]\textsf{\textbf{\color{ink}C.}} #3\\[3.5pt]\textsf{\textbf{\color{ink}D.}} #4\end{tabularx}}
\newcommand{\answerbox}{\hfill\mbox{\textsf{\footnotesize\textit{\color{mutedSlate}Đáp số:}}\ \tikz[baseline=-3pt]\node[draw=brand,line width=0.65pt,fill=brandTint,rounded corners=2.5pt,minimum width=1.85cm,minimum height=13.5pt,inner sep=0pt]{};}}
\newcommand{\baibadge}[2]{\par\vspace{2.5pt}\needspace{3\baselineskip}\tikz[baseline=(b.base)]\node[fill=brandDark,text=white,rounded corners=2pt,inner xsep=5pt,inner ysep=1.5pt,font=\small\sffamily\bfseries] (b) {#1};\ \textsf{\footnotesize\textbf{\color{amberAccent}#2}}\ }

\setlist[enumerate,1]{label=\textsf{\textbf{\color{brandDark}Câu \arabic*.}},leftmargin=1.18cm,labelsep=.14cm,itemsep=2.2pt,topsep=1.8pt,parsep=0pt}
\setlist[enumerate,2]{label=\textsf{\textbf{\color{slateText}\alph*)}},leftmargin=.55cm,itemsep=1.5pt,topsep=0.5pt,parsep=0pt}
\newcolumntype{C}[1]{>{\centering\arraybackslash}m{#1}}


\begin{document}

\begin{tcolorbox}[enhanced,colback=brandTint,colframe=brandLine,boxrule=0.7pt,arc=4pt,borderline west={3.5pt}{0pt}{brandDark},left=9pt,right=9pt,top=6pt,bottom=6pt,before skip=0pt,after skip=4pt]
\noindent
\begin{minipage}[c]{0.34\textwidth}
    \centering
    {\footnotesize\sffamily\bfseries\color{brandDark} LÀM TOÁN ĐỂ HẠNH PHÚC}\\[2pt]
    {\color{brandLine}\rule{3.0cm}{0.8pt}}\\[3pt]
    {\large\sffamily\bfseries\color{ink} ĐỀ ÔN TẬP SỐ 6}
\end{minipage}\hfill
\begin{minipage}[c]{0.49\textwidth}
    \centering
    {\normalsize\sffamily\bfseries\color{ink} ĐỀ ÔN TẬP TOÁN 10 — CHƯƠNG IV. VECTƠ}\\[2pt]
    {\small\sffamily\bfseries\color{slateText} NĂM HỌC 2026 -- 2027}\\[1pt]
    {\footnotesize\sffamily\textbf{\color{brandDark}Môn: TOÁN -- Khối lớp 10}}\\
    {\footnotesize\sffamily\textit{\color{mutedSlate}Thời gian làm bài: 90 phút (không kể thời gian phát đề)}}
\end{minipage}\hfill
\begin{minipage}[c]{0.15\textwidth}
    \raggedleft
    \begin{tcolorbox}[colback=amberTint,colframe=amberAccent,boxrule=0.8pt,arc=3pt,left=3pt,right=3pt,top=3pt,bottom=3pt,boxsep=1pt]
        \centering
        {\scriptsize\sffamily\bfseries\color{amberAccent} MÃ ĐỀ}\\[-1pt]
        {\normalsize\sffamily\bfseries\color{ink} 1047}
    \end{tcolorbox}
\end{minipage}
\end{tcolorbox}

\begin{tcolorbox}[colback=white,colframe=ruleGray,boxrule=0.5pt,arc=3pt,left=8pt,right=8pt,top=3.5pt,bottom=3.5pt,before skip=0pt,after skip=4pt]
\noindent\textsf{\small\textbf{\color{ink}Họ và tên thí sinh:}}\dotfill\quad\textsf{\small\textbf{\color{ink}SBD:}}\makebox[3.2cm]{\dotfill}\quad\textsf{\small\textbf{\color{ink}Lớp:}}\makebox[1.8cm]{\dotfill}
\end{tcolorbox}

\parttitle{PHẦN I. Câu trắc nghiệm nhiều phương án lựa chọn}
\textit{\color{mutedSlate}\small Thí sinh trả lời từ câu 1 đến câu 20. Mỗi câu hỏi thí sinh chỉ chọn một phương án.}
\begin{enumerate}
\item Cho $6$ điểm phân biệt nằm trên cùng một đường thẳng. Số vectơ khác vectơ-không có điểm đầu và điểm cuối là $2$ trong $6$ điểm đã cho là
\four{$15$.}{$30$.}{$12$.}{$36$.}

\item Cho hai vectơ $\vec{a}$ và $\vec{b}$ khác vectơ-không và ngược hướng với nhau, biết $|\vec{a}| = 6$ và $|\vec{b}| = 2$. Khẳng định nào sau đây đúng?
\four{$\vec{a} = 3\vec{b}$.}{$\vec{a} = -\dfrac{1}{3}\vec{b}$.}{$\vec{b} = 3\vec{a}$.}{$\vec{a} = -3\vec{b}$.}

\item Cho lục giác đều $ABCDEF$ tâm $O$. Tổng vectơ $\vec{OA} + \vec{OB} + \vec{OC} + \vec{OD} + \vec{OE} + \vec{OF}$ bằng
\four{$\vec{0}$.}{$2\vec{AD}$.}{$6\vec{OA}$.}{$\vec{AB}$.}

\item Cho hình chữ nhật $ABCD$ có $AB = 5$ và $AD = 12$. Độ dài của vectơ $\vec{BA} + \vec{BC}$ bằng
\four{$17$.}{$7$.}{$13$.}{$\sqrt{119}$.}

\item Cho hình vuông $ABCD$ có cạnh bằng $2a$ và tâm $O$. Độ dài của vectơ $\vec{OA} + \vec{OB}$ bằng
\four{$a$.}{$a\sqrt{2}$.}{$2a\sqrt{2}$.}{$2a$.}

\item Cho tam giác $ABC$. Gọi $I$ và $K$ lần lượt là trung điểm của các cạnh $AB$ và $AC$. Mệnh đề nào sau đây đúng?
\four{$\vec{IK} = 2\vec{BC}$.}{$\vec{IK} = \dfrac{1}{2}\vec{BC}$.}{$\vec{IK} = -\dfrac{1}{2}\vec{BC}$.}{$\vec{IK} = \dfrac{1}{3}\vec{BC}$.}

\item Cho tam giác $ABC$ và điểm $M$ thỏa mãn $\vec{MA} + 2\vec{MB} = \vec{0}$. Đẳng thức nào sau đây đúng?
\two{$\vec{CM} = \dfrac{1}{3}\vec{CA} + \dfrac{2}{3}\vec{CB}$.}{$\vec{CM} = \dfrac{2}{3}\vec{CA} + \dfrac{1}{3}\vec{CB}$.}{$\vec{CM} = \dfrac{1}{3}\vec{CA} - \dfrac{2}{3}\vec{CB}$.}{$\vec{CM} = -\dfrac{1}{3}\vec{CA} + \dfrac{2}{3}\vec{CB}$.}

\item Cho hình bình hành $ABCD$ tâm $O$. Điểm $M$ thỏa mãn $\vec{MA} + \vec{MB} + \vec{MC} + \vec{MD} = \vec{AB}$. Khẳng định nào sau đây đúng?
\four{$\vec{OM} = \dfrac{1}{4}\vec{AB}$.}{$\vec{MO} = \dfrac{1}{2}\vec{AB}$.}{$\vec{MO} = \dfrac{1}{4}\vec{AB}$.}{$\vec{MO} = -\dfrac{1}{4}\vec{AB}$.}

\item Trong mặt phẳng tọa độ $Oxy$, cho hai vectơ $\vec{a} = (3; -4)$ và $\vec{b} = (-1; 2)$. Tọa độ của vectơ $\vec{c} = 2\vec{a} + 3\vec{b}$ là
\four{$(9; -14)$.}{$(5; -2)$.}{$(-3; 2)$.}{$(3; -2)$.}

\item Trong mặt phẳng tọa độ $Oxy$, cho điểm $A(4; -2)$ và vectơ $\vec{AB} = (-3; 5)$. Tọa độ của điểm $B$ là
\four{$B(1; 3)$.}{$B(7; -7)$.}{$B(-1; -3)$.}{$B(-7; 7)$.}

\item Trong mặt phẳng tọa độ $Oxy$, cho tam giác $ABC$ có $A(-2; 5)$, $B(4; -1)$ và trọng tâm $G(1; 2)$. Tọa độ đỉnh $C$ là
\four{$C(3; 6)$.}{$C(-1; 2)$.}{$C(1; 2)$.}{$C(1; 6)$.}

\item Trong mặt phẳng tọa độ $Oxy$, cho ba điểm $P(-1; 3)$, $Q(2; -3)$ và $R(5; m)$. Giá trị của tham số $m$ để ba điểm $P, Q, R$ thẳng hàng là
\four{$m = 9$.}{$m = -9$.}{$m = -3$.}{$m = 15$.}
\pagebreak

\item Trong mặt phẳng tọa độ $Oxy$, cho hai vectơ $\vec{u} = (m; -6)$ và $\vec{v} = (-3; m + 3)$. Tất cả các giá trị của $m$ để $\vec{u}$ cùng phương với $\vec{v}$ là
\four{$m = 3$.}{$m = -3$ hoặc $m = 6$.}{$m = 3$ hoặc $m = -6$.}{$m = -6$.}

\item Cho tam giác $ABC$ cân tại $A$ có $\widehat{A} = 50^\circ$. Góc giữa hai vectơ $\vec{AB}$ và $\vec{BC}$ bằng
\four{$115^\circ$.}{$65^\circ$.}{$50^\circ$.}{$130^\circ$.}

\item Cho hai vectơ $\vec{a}, \vec{b}$ thỏa mãn $|\vec{a}| = 4$, $|\vec{b}| = 5$ và $(\vec{a}, \vec{b}) = 120^\circ$. Giá trị của $\vec{a} \cdot (\vec{a} + \vec{b})$ bằng
\four{$26$.}{$16$.}{$-10$.}{$6$.}

\item Trong mặt phẳng tọa độ $Oxy$, cho tam giác $ABC$ có $A(1; 2)$, $B(4; -2)$ và $C(-1; 4)$. Giá trị của tích vô hướng $\vec{AB} \cdot \vec{AC}$ bằng
\four{$14$.}{$-14$.}{$2$.}{$-2$.}

\item Trong mặt phẳng tọa độ $Oxy$, cho vectơ $\vec{u} = (x; x + 2)$ với $x > 0$. Biết $|\vec{u}| = 10$, giá trị của $x$ bằng
\four{$6$.}{$8$.}{$5$.}{$4$.}

\item Trong mặt phẳng tọa độ $Oxy$, cho hai vectơ $\vec{a} = (2; -5)$ và $\vec{b} = (m + 4; 2)$. Giá trị của $m$ để $\vec{a} \perp \vec{b}$ là
\four{$m = -1$.}{$m = 5$.}{$m = -9$.}{$m = 1$.}

\item Trong mặt phẳng tọa độ $Oxy$, cho hai vectơ $\vec{a} = (3; 1)$ và $\vec{b} = (1; 2)$. Giá trị của $\cos(\vec{a}, \vec{b})$ bằng
\four{$\dfrac{1}{2}$.}{$\dfrac{\sqrt{3}}{2}$.}{$\dfrac{\sqrt{2}}{2}$.}{$\dfrac{1}{\sqrt{5}}$.}

\item Hai lực $\vec{F}_1$ và $\vec{F}_2$ có cùng điểm đặt, cùng độ lớn bằng $30\text{ N}$ và hợp với nhau một góc $60^\circ$. Độ lớn của hợp lực $\vec{F} = \vec{F}_1 + \vec{F}_2$ bằng
\four{$30\text{ N}$.}{$30\sqrt{3}\text{ N}$.}{$60\text{ N}$.}{$30\sqrt{2}\text{ N}$.}
\end{enumerate}

\parttitle{PHẦN II. Câu trắc nghiệm đúng sai}
\textit{\color{mutedSlate}\small Thí sinh trả lời từ câu 1 đến câu 4. Trong mỗi ý a), b), c), d) ở mỗi câu, thí sinh chọn đúng hoặc sai.}
\begin{enumerate}
\item Cho hình chữ nhật $ABCD$ có $AB = 8$, $AD = 12$ và $M$ là trung điểm của cạnh $BC$ (xem hình vẽ bên).
\par\noindent\begin{minipage}[t]{0.73\linewidth}
\begin{enumerate}
    \item Độ dài đường chéo $AC = 20$.
    \item Tích vô hướng $\vec{AB} \cdot \vec{AC} = 64$.
    \item Tích vô hướng $\vec{AM} \cdot \vec{BD} = 16$.
    \item Độ dài đoạn thẳng $AM = 10$.
\end{enumerate}
\end{minipage}\hfill
\begin{minipage}[t]{0.25\linewidth}\raggedleft\vspace{-10pt}
\begin{tikzpicture}[scale=0.62, >=Stealth]
  \coordinate (A) at (0,0);
  \coordinate (B) at (2.0,0);
  \coordinate (C) at (2.0,2.8);
  \coordinate (D) at (0,2.8);
  \coordinate (M) at (2.0,1.4);
  \fill[brandTint!75] (A) -- (B) -- (C) -- (D) -- cycle;
  \draw[thick, ink] (A) -- (B) -- (C) -- (D) -- cycle;
  \draw[dashed, mutedSlate] (A) -- (C);
  \draw[dashed, amberAccent, thick] (A) -- (M);
  \draw[brand] (0,0.28) -- (0.28,0.28) -- (0.28,0);
  \fill[amberAccent] (M) circle (2.2pt);
  \node[below left, font=\scriptsize\sffamily] at (A) {$A$};
  \node[below right, font=\scriptsize\sffamily] at (B) {$B$};
  \node[above right, font=\scriptsize\sffamily] at (C) {$C$};
  \node[above left, font=\scriptsize\sffamily] at (D) {$D$};
  \node[right, font=\scriptsize\sffamily] at (M) {$M$};
\end{tikzpicture}
\end{minipage}

\item Cho tam giác $ABC$. Gọi $G$ là trọng tâm của tam giác $ABC$, $M$ là trung điểm của cạnh $BC$, và điểm $N$ thỏa mãn $\vec{AN} = 3\vec{AB}$.
\begin{enumerate}
    \item $\vec{GA} + \vec{GB} + \vec{GC} = \vec{0}$.
    \item $\vec{MG} = \dfrac{1}{6}\vec{AB} + \dfrac{1}{6}\vec{AC}$.
    \item $\vec{MN} = \dfrac{5}{2}\vec{AB} - \dfrac{1}{2}\vec{AC}$.
    \item $\vec{CN} = 3\vec{AB} + \vec{AC}$.
\end{enumerate}

\item Trong mặt phẳng tọa độ $Oxy$, cho tam giác $ABC$ có $A(-1; 2)$, $B(3; 4)$ và $C(1; -2)$.
\begin{enumerate}
    \item Tọa độ vectơ $\vec{AB} = (-4; -2)$.
    \item Trọng tâm $G$ của tam giác $ABC$ có tọa độ là $G\!\left(1; \dfrac{2}{3}\right)$.
    \item Tam giác $ABC$ vuông cân tại $A$.
    \item Diện tích tam giác $ABC$ bằng $10$.
\end{enumerate}

\item Cho tam giác $ABC$ có $AB = 6$, $AC = 10$ và $\widehat{BAC} = 120^\circ$.
\begin{enumerate}
    \item Tích vô hướng $\vec{AB} \cdot \vec{AC} = -30$.
    \item Độ dài cạnh $BC = 14$.
    \item Độ dài vectơ $\vec{AB} + \vec{AC}$ bằng $16$.
    \item Tích vô hướng $\vec{BA} \cdot \vec{BC} = 6$.
\end{enumerate}
\end{enumerate}
\pagebreak

\parttitle{PHẦN III. Câu trắc nghiệm trả lời ngắn}
\textit{\color{mutedSlate}\small Thí sinh trả lời từ câu 1 đến câu 6.}
\begin{enumerate}
\item Cho hình thoi $ABCD$ có cạnh bằng $12$ và $\widehat{BAD} = 60^\circ$. Tính độ dài của vectơ $\vec{u} = \vec{AB} - \vec{AD}$. \answerbox

\item Cho tam giác $ABC$. Gọi $M$ là trung điểm của cạnh $AB$ và điểm $N$ nằm trên cạnh $AC$ sao cho $AN = 3NC$. Biết rằng $\vec{MN} = x\vec{AB} + y\vec{AC}$. Tính giá trị biểu thức $T = 8y - 2x$. \answerbox

\item Trong mặt phẳng tọa độ $Oxy$, cho ba điểm $A(2; -1)$, $B(0; 3)$ và $C(-3; 2)$. Điểm $M(x_M; y_M)$ thỏa mãn đẳng thức $\vec{MA} + 3\vec{MB} - 2\vec{MC} = \vec{0}$. Tính giá trị $T = 2x_M + y_M$. \answerbox

\item Trong mặt phẳng tọa độ $Oxy$, cho hai vectơ $\vec{u} = (4; -3)$ và $\vec{v} = (8; 6)$. Tính $\cos(\vec{u}, \vec{v})$ (viết kết quả dưới dạng số thập phân). \answerbox

\item Trong mặt phẳng tọa độ $Oxy$, cho hai điểm $A(-1; 2)$, $B(5; 4)$ và điểm $M(m; 0)$. Tìm giá trị của $m$ sao cho $MA = MB$. \answerbox

\item Hai tàu kéo một xà lan trên mặt nước bằng hai lực $\vec{F}_1$ và $\vec{F}_2$ cùng nằm trong mặt phẳng nằm ngang, có độ lớn lần lượt là $F_1 = 300\text{ N}$ và $F_2 = 500\text{ N}$, hợp với nhau một góc $60^\circ$ (như hình vẽ). Tính độ lớn của hợp lực $\vec{F} = \vec{F}_1 + \vec{F}_2$ (đơn vị: Newton).
\par\noindent\begin{minipage}[c]{0.68\linewidth}
\answerbox
\end{minipage}\hfill
\begin{minipage}[c]{0.30\linewidth}\raggedleft
\begin{tikzpicture}[scale=0.56, >=Stealth]
  \fill[brandTint] (-0.35,-0.25) rectangle (0.35,0.25);
  \draw[thick, ink, rounded corners=1.5pt] (-0.35,-0.25) rectangle (0.35,0.25);
  \coordinate (O) at (0.35,0);
  \coordinate (F1) at (2.35,0);
  \coordinate (F2) at (1.95, 2.77);
  \draw[->, very thick, brandDark] (O) -- (F1) node[below, font=\scriptsize] {$\vec{F}_1$};
  \draw[->, very thick, amberAccent] (O) -- (F2) node[above left, font=\scriptsize] {$\vec{F}_2$};
  \pic[draw=mutedSlate, angle radius=0.45cm, "$60^\circ$", angle eccentricity=1.65, font=\tiny] {angle = F1--O--F2};
\end{tikzpicture}
\end{minipage}
\end{enumerate}

\parttitle{PHẦN IV. Câu tự luận}
\textit{\color{mutedSlate}\small Thí sinh trình bày lời giải chi tiết từ Bài 1 đến Bài 3.}

\baibadge{Bài 1}{(1,0 điểm)} Cho tam giác $ABC$. Gọi $I$ là trung điểm của cạnh $BC$, $M$ là trung điểm của đoạn thẳng $AI$, và điểm $N$ thuộc cạnh $AC$ sao cho $NC = 2NA$.
\begin{enumerate}[label=\alph*), leftmargin=0.6cm, itemsep=2pt]
    \item (0,5 điểm) Hãy biểu diễn hai vectơ $\vec{BM}$ và $\vec{BN}$ theo hai vectơ $\vec{AB}$ và $\vec{AC}$.
    \item (0,5 điểm) Chứng minh rằng ba điểm $B, M, N$ thẳng hàng.
\end{enumerate}

\baibadge{Bài 2}{(1,0 điểm)} Trong mặt phẳng tọa độ $Oxy$, cho ba điểm $A(2; 5)$, $B(-1; -1)$ và $C(4; 4)$.
\begin{enumerate}[label=\alph*), leftmargin=0.6cm, itemsep=2pt]
    \item (0,5 điểm) Chứng minh tam giác $ABC$ vuông tại $A$ và tìm tọa độ điểm $D$ sao cho tứ giác $ABCD$ là hình bình hành.
    \item (0,5 điểm) Tìm tọa độ điểm $P$ nằm trên trục hoành $Ox$ sao cho biểu thức $T = |\vec{PA} + \vec{PB} + 2\vec{PC}|$ đạt giá trị nhỏ nhất, và tính giá trị nhỏ nhất đó.
\end{enumerate}

\baibadge{Bài 3}{(1,0 điểm)} Một chiếc phà muốn đi thẳng qua một con sông rộng $240\text{ m}$ (vuông góc với hai bờ sông song song) từ bến $A$ ở bờ Nam sang bến $B$ ở bờ Bắc. Biết nước sông chảy theo hướng từ Tây sang Đông với vận tốc $\vec{v}_n$ có độ lớn $v_n = 3\text{ m/s}$, còn vận tốc thực của phà khi nước đứng yên có độ lớn $v_p = 5\text{ m/s}$ theo vectơ $\vec{v}_p$. Chọn hệ trục tọa độ $Oxy$ gắn với mặt đất có gốc $O \equiv A$, tia $Ox$ hướng về phía Đông và tia $Oy$ hướng về phía Bắc (đơn vị trên trục vận tốc là $\text{m/s}$, xem hình vẽ).
\par\noindent\begin{minipage}[t]{0.66\linewidth}
\begin{enumerate}[label=\alph*), leftmargin=0.6cm, itemsep=2pt]
    \item (0,5 điểm) Để phà đi thẳng theo hướng $AB$ (trùng tia $Oy$), vận tốc tổng hợp $\vec{v} = \vec{v}_p + \vec{v}_n$ phải cùng hướng với tia $Oy$. Hãy tìm tọa độ của vectơ $\vec{v}_p$ và vectơ $\vec{v}$ trong hệ trục $Oxy$.
    \item (0,5 điểm) Tính thời gian để phà sang đến bến $B$ và xác định góc hợp bởi hướng mũi phà (vectơ $\vec{v}_p$) với hướng chảy của dòng nước (tia $Ox$, làm tròn kết quả góc đến hàng phần mười của độ).
\end{enumerate}
\end{minipage}\hfill
\begin{minipage}[t]{0.32\linewidth}\raggedleft\vspace{-10pt}
\begin{tikzpicture}[scale=0.62, >=Stealth]
  \fill[brandTint] (-1.8,0) rectangle (2.3,2.0);
  \draw[thick, brandDark] (-1.8,0) -- (2.3,0);
  \draw[thick, brandDark] (-1.8,2.0) -- (2.3,2.0);
  \coordinate (O) at (0,0);
  \coordinate (B) at (0,2.0);
  \coordinate (Vn) at (1.125,0);
  \coordinate (Vp) at (-1.125,1.5);
  \coordinate (V) at (0,1.5);
  \draw[->, mutedSlate] (O) -- (2.25,0) node[below, font=\tiny\sffamily] {$x$ (Đông)};
  \draw[->, mutedSlate] (O) -- (0,2.55) node[above, font=\tiny\sffamily] {$y$ (Bắc)};
  \draw[->, very thick, amberAccent] (O) -- (Vn) node[below, font=\tiny] {$\vec{v}_n$};
  \draw[->, very thick, ink] (O) -- (Vp) node[midway, below left, font=\tiny] {$\vec{v}_p$};
  \draw[->, very thick, brand] (O) -- (V) node[right, font=\tiny] {$\vec{v}$};
  \draw[dashed, mutedSlate!80] (Vp) -- (V) -- (Vn);
  \draw[<->, mutedSlate] (2.05,0) -- (2.05,2.0) node[midway, right, font=\tiny] {$240\text{ m}$};
  \fill[ink] (O) circle (1.6pt);
  \fill[ink] (B) circle (1.6pt);
  \node[below left, font=\tiny] at (O) {$A \equiv O$};
  \node[above left, font=\tiny] at (B) {$B$};
\end{tikzpicture}
\end{minipage}


\newpage
\pagestyle{dapan}
\setcounter{page}{1}

\begin{tcolorbox}[enhanced,colback=brandTint,colframe=brandLine,boxrule=0.7pt,arc=4pt,borderline west={3.5pt}{0pt}{brandDark},top=5pt,bottom=5pt,before skip=0pt,after skip=5pt]
  \centering
  {\Large\sffamily\bfseries\color{ink} BẢNG ĐÁP ÁN TỔNG HỢP}\\[2pt]
  {\small\sffamily\color{mutedSlate} Đề ôn tập Toán 10 — Chương IV. Vectơ — Đề số 6 (Mã đề 1047)}
\end{tcolorbox}

\noindent\textsf{\small\textbf{\color{ink}PHẦN I. Câu trắc nghiệm nhiều phương án lựa chọn (20 câu)}}
\vspace{2pt}

\begin{center}
\renewcommand{\arraystretch}{1.22}
\arrayrulecolor{ruleGray}
\begin{tabular}{|C{0.95cm}|C{0.95cm}|C{0.95cm}|C{0.95cm}|C{0.95cm}|C{0.95cm}|C{0.95cm}|C{0.95cm}|C{0.95cm}|C{0.95cm}|}
\hline
\rowcolor{brandDark}
\color{white}\textsf{\textbf{1}} & \color{white}\textsf{\textbf{2}} & \color{white}\textsf{\textbf{3}} & \color{white}\textsf{\textbf{4}} & \color{white}\textsf{\textbf{5}} & \color{white}\textsf{\textbf{6}} & \color{white}\textsf{\textbf{7}} & \color{white}\textsf{\textbf{8}} & \color{white}\textsf{\textbf{9}} & \color{white}\textsf{\textbf{10}} \\ \hline
\rowcolor{brandTint}
\textbf{\color{ink}B} & \textbf{\color{ink}D} & \textbf{\color{ink}A} & \textbf{\color{ink}C} & \textbf{\color{ink}D} & \textbf{\color{ink}B} & \textbf{\color{ink}A} & \textbf{\color{ink}C} & \textbf{\color{ink}D} & \textbf{\color{ink}A} \\ \hline
\rowcolor{brandDark}
\color{white}\textsf{\textbf{11}} & \color{white}\textsf{\textbf{12}} & \color{white}\textsf{\textbf{13}} & \color{white}\textsf{\textbf{14}} & \color{white}\textsf{\textbf{15}} & \color{white}\textsf{\textbf{16}} & \color{white}\textsf{\textbf{17}} & \color{white}\textsf{\textbf{18}} & \color{white}\textsf{\textbf{19}} & \color{white}\textsf{\textbf{20}} \\ \hline
\rowcolor{brandTint}
\textbf{\color{ink}C} & \textbf{\color{ink}B} & \textbf{\color{ink}C} & \textbf{\color{ink}A} & \textbf{\color{ink}D} & \textbf{\color{ink}B} & \textbf{\color{ink}A} & \textbf{\color{ink}D} & \textbf{\color{ink}C} & \textbf{\color{ink}B} \\ \hline
\end{tabular}
\end{center}

\vspace{2pt}
\noindent\begin{minipage}[t]{0.55\linewidth}
\textsf{\small\textbf{\color{ink}PHẦN II. Câu trắc nghiệm đúng/sai}}\\[2pt]
\renewcommand{\arraystretch}{1.22}
\arrayrulecolor{ruleGray}
\begin{tabular}{|C{0.9cm}|C{0.7cm}|C{0.7cm}|C{0.7cm}|C{0.7cm}|}
\hline
\rowcolor{brandDark}
\color{white}\textsf{\textbf{Câu}} & \color{white}\textsf{\textbf{a)}} & \color{white}\textsf{\textbf{b)}} & \color{white}\textsf{\textbf{c)}} & \color{white}\textsf{\textbf{d)}} \\ \hline
\rowcolor{brandTint}
\textbf{\color{ink}1} & \color{mutedSlate}\textbf{S} & \color{brandDark}\textbf{Đ} & \color{mutedSlate}\textbf{S} & \color{brandDark}\textbf{Đ} \\ \hline
\textbf{\color{ink}2} & \color{brandDark}\textbf{Đ} & \color{mutedSlate}\textbf{S} & \color{brandDark}\textbf{Đ} & \color{mutedSlate}\textbf{S} \\ \hline
\rowcolor{brandTint}
\textbf{\color{ink}3} & \color{mutedSlate}\textbf{S} & \color{mutedSlate}\textbf{S} & \color{brandDark}\textbf{Đ} & \color{brandDark}\textbf{Đ} \\ \hline
\textbf{\color{ink}4} & \color{brandDark}\textbf{Đ} & \color{brandDark}\textbf{Đ} & \color{mutedSlate}\textbf{S} & \color{mutedSlate}\textbf{S} \\ \hline
\end{tabular}
\end{minipage}\hfill
\begin{minipage}[t]{0.42\linewidth}
\textsf{\small\textbf{\color{ink}PHẦN III. Trả lời ngắn}}\\[2pt]
\renewcommand{\arraystretch}{1.22}
\arrayrulecolor{ruleGray}
\begin{tabular}{|C{0.85cm}|C{1.35cm}|C{0.85cm}|C{1.35cm}|}
\hline
\rowcolor{brandDark}
\color{white}\textsf{\textbf{Câu}} & \color{white}\textsf{\textbf{Đáp số}} & \color{white}\textsf{\textbf{Câu}} & \color{white}\textsf{\textbf{Đáp số}} \\ \hline
\rowcolor{brandTint}
\textbf{\color{ink}1} & \textbf{$12$} & \textbf{\color{ink}4} & \textbf{$0{,}28$} \\ \hline
\textbf{\color{ink}2} & \textbf{$7$}  & \textbf{\color{ink}5} & \textbf{$3$} \\ \hline
\rowcolor{brandTint}
\textbf{\color{ink}3} & \textbf{$10$} & \textbf{\color{ink}6} & \textbf{$700$} \\ \hline
\end{tabular}
\end{minipage}

\vspace{4pt}
\begin{tcolorbox}[enhanced,colback=amberTint,colframe=ruleGray,boxrule=0.5pt,arc=3pt,borderline west={3pt}{0pt}{amberAccent},left=7pt,right=7pt,top=4pt,bottom=4pt,before skip=2pt,after skip=4pt]
\noindent\textsf{\small\textbf{\color{ink}PHẦN IV. Đáp số tự luận:}}\\[2pt]
\small
\textbullet\ \textbf{\color{ink}Bài 1:} a) $\vec{BM} = -\dfrac{3}{4}\vec{AB} + \dfrac{1}{4}\vec{AC}$; $\vec{BN} = -\vec{AB} + \dfrac{1}{3}\vec{AC}$. b) $\vec{BM} = \dfrac{3}{4}\vec{BN} \Rightarrow B, M, N$ thẳng hàng.\\
\textbullet\ \textbf{\color{ink}Bài 2:} a) $\vec{AB}\cdot\vec{AC} = 0$; $D(7; 10)$. b) $P\!\left(\dfrac{9}{4}; 0\right)$, GTNN bằng $12$.\\
\textbullet\ \textbf{\color{ink}Bài 3:} a) $\vec{v}_p = (-3; 4)\text{ m/s}$, $\vec{v} = (0; 4)\text{ m/s}$. b) $t = 60\text{ s}$; góc $\approx 126{,}9^\circ$.
\end{tcolorbox}

\parttitle{\centering HƯỚNG DẪN GIẢI CHI TIẾT}
\vspace{1pt}

\begin{multicols}{2}
\raggedcolumns
\small
\arrayrulecolor{black}

\textbf{\color{accent}\textsf{PHẦN I. TRẮC NGHIỆM}}

\textbf{Câu 1.} Số vectơ khác $\vec{0}$ tạo bởi $2$ trong $6$ điểm là $A_6^2 = 6 \times 5 = 30$. Chọn \textbf{B}.

\textbf{Câu 2.} Vì $\vec{a}, \vec{b}$ ngược hướng và $|\vec{a}| = 3|\vec{b}|$ nên $\vec{a} = -3\vec{b}$. Chọn \textbf{D}.

\textbf{Câu 3.} Nhóm các cặp vectơ đối qua tâm $O$: $(\vec{OA} + \vec{OD}) + (\vec{OB} + \vec{OE}) + (\vec{OC} + \vec{OF}) = \vec{0}$. Chọn \textbf{A}.

\textbf{Câu 4.} $\vec{BA} + \vec{BC} = \vec{BD} \Rightarrow |\vec{BA} + \vec{BC}| = BD = \sqrt{5^2 + 12^2} = 13$. Chọn \textbf{C}.

\textbf{Câu 5.} Gọi $M$ là trung điểm $AB \Rightarrow |\vec{OA} + \vec{OB}| = 2OM = 2a$. Chọn \textbf{D}.

\textbf{Câu 6.} $IK$ là đường trung bình $\triangle ABC$ nên $\vec{IK} = \dfrac{1}{2}\vec{BC}$. Chọn \textbf{B}.

\textbf{Câu 7.} $\vec{MA} + 2\vec{MB} = \vec{0} \Leftrightarrow \vec{CM} = \dfrac{1}{3}\vec{CA} + \dfrac{2}{3}\vec{CB}$. Chọn \textbf{A}.

\textbf{Câu 8.} $O$ là tâm $ABCD \Rightarrow \vec{MA} + \vec{MB} + \vec{MC} + \vec{MD} = 4\vec{MO} = \vec{AB} \Rightarrow \vec{MO} = \dfrac{1}{4}\vec{AB}$. Chọn \textbf{C}.

\textbf{Câu 9.} $\vec{c} = 2(3; -4) + 3(-1; 2) = (6 - 3; -8 + 6) = (3; -2)$. Chọn \textbf{D}.

\textbf{Câu 10.} $B = A + \vec{AB} = (4 - 3; -2 + 5) = (1; 3)$. Chọn \textbf{A}.

\textbf{Câu 11.} $C = 3G - A - B = (3 - (-2) - 4; 6 - 5 - (-1)) = (1; 2)$. Chọn \textbf{C}.

\textbf{Câu 12.} $\vec{PQ} = (3; -6)$, $\vec{PR} = (6; m - 3)$. Thẳng hàng $\Leftrightarrow 3(m - 3) + 36 = 0 \Leftrightarrow m = -9$. Chọn \textbf{B}.

\textbf{Câu 13.} $\vec{u} \parallel \vec{v} \Leftrightarrow m(m + 3) - 18 = 0 \Leftrightarrow m^2 + 3m - 18 = 0 \Leftrightarrow m \in \{3; -6\}$. Chọn \textbf{C}.

\textbf{Câu 14.} $\triangle ABC$ cân tại $A$ có $\widehat{B} = \dfrac{180^\circ - 50^\circ}{2} = 65^\circ \Rightarrow (\vec{AB}, \vec{BC}) = 180^\circ - 65^\circ = 115^\circ$. Chọn \textbf{A}.

\textbf{Câu 15.} $\vec{a}\cdot(\vec{a} + \vec{b}) = |\vec{a}|^2 + |\vec{a}||\vec{b}|\cos 120^\circ = 16 + 4\cdot 5\left(-\dfrac{1}{2}\right) = 6$. Chọn \textbf{D}.

\textbf{Câu 16.} $\vec{AB} = (3; -4)$, $\vec{AC} = (-2; 2) \Rightarrow \vec{AB}\cdot\vec{AC} = -6 - 8 = -14$. Chọn \textbf{B}.

\textbf{Câu 17.} $|\vec{u}| = 10 \Leftrightarrow x^2 + (x + 2)^2 = 100 \Leftrightarrow x^2 + 2x - 48 = 0 \Leftrightarrow x = 6$ (do $x > 0$). Chọn \textbf{A}.

\textbf{Câu 18.} $\vec{a} \perp \vec{b} \Leftrightarrow 2(m + 4) - 10 = 0 \Leftrightarrow 2m - 2 = 0 \Leftrightarrow m = 1$. Chọn \textbf{D}.

\textbf{Câu 19.} $\cos(\vec{a}, \vec{b}) = \dfrac{3\cdot 1 + 1\cdot 2}{\sqrt{10}\cdot\sqrt{5}} = \dfrac{5}{5\sqrt{2}} = \dfrac{\sqrt{2}}{2}$. Chọn \textbf{C}.

\textbf{Câu 20.} $|\vec{F}| = \sqrt{30^2 + 30^2 + 2\cdot 30\cdot 30\cos 60^\circ} = 30\sqrt{3}\text{ N}$. Chọn \textbf{B}.

\textbf{\color{brandDark}\textsf{PHẦN II. ĐÚNG / SAI}}

\textbf{Câu 1.}
a) $AC = \sqrt{8^2 + 12^2} = 4\sqrt{13} \neq 20$ (\textbf{S}).
b) $\vec{AB}\cdot\vec{AC} = \vec{AB}\cdot(\vec{AB} + \vec{AD}) = AB^2 + 0 = 64$ (\textbf{Đ}).
c) $\vec{AM} = \vec{AB} + \dfrac{1}{2}\vec{AD}$, $\vec{BD} = \vec{AD} - \vec{AB} \Rightarrow \vec{AM}\cdot\vec{BD} = -64 + 72 = 8 \neq 16$ (\textbf{S}).
d) $AM = \sqrt{AB^2 + BM^2} = \sqrt{8^2 + 6^2} = 10$ (\textbf{Đ}).

\textbf{Câu 2.}
a) Theo tính chất trọng tâm, $\vec{GA} + \vec{GB} + \vec{GC} = \vec{0}$ (\textbf{Đ}).
b) $\vec{MG} = \vec{AG} - \vec{AM} = \left(\dfrac{1}{3} - \dfrac{1}{2}\right)(\vec{AB} + \vec{AC}) = -\dfrac{1}{6}\vec{AB} - \dfrac{1}{6}\vec{AC}$ (\textbf{S}).
c) $\vec{MN} = \vec{AN} - \vec{AM} = 3\vec{AB} - \dfrac{1}{2}\vec{AB} - \dfrac{1}{2}\vec{AC} = \dfrac{5}{2}\vec{AB} - \dfrac{1}{2}\vec{AC}$ (\textbf{Đ}).
d) $\vec{CN} = \vec{AN} - \vec{AC} = 3\vec{AB} - \vec{AC}$ (\textbf{S}).

\textbf{Câu 3.}
a) $\vec{AB} = (3 - (-1); 4 - 2) = (4; 2) \neq (-4; -2)$ (\textbf{S}).
b) $G = \left(\dfrac{-1 + 3 + 1}{3}; \dfrac{2 + 4 - 2}{3}\right) = \left(1; \dfrac{4}{3}\right)$ (\textbf{S}).
c) $\vec{AB} = (4; 2)$, $\vec{AC} = (2; -4) \Rightarrow \vec{AB}\cdot\vec{AC} = 0$ và $AB = AC = 2\sqrt{5}$ (\textbf{Đ}).
d) $S_{\triangle ABC} = \dfrac{1}{2}AB\cdot AC = \dfrac{1}{2}(2\sqrt{5})^2 = 10$ (\textbf{Đ}).

\textbf{Câu 4.}
a) $\vec{AB}\cdot\vec{AC} = 6\cdot 10\cdot\cos 120^\circ = -30$ (\textbf{Đ}).
b) $BC = \sqrt{6^2 + 10^2 - 2(-30)} = \sqrt{196} = 14$ (\textbf{Đ}).
c) $|\vec{AB} + \vec{AC}| = \sqrt{36 + 100 + 2(-30)} = \sqrt{76} = 2\sqrt{19} \neq 16$ (\textbf{S}).
d) $\vec{BA}\cdot\vec{BC} = (-\vec{AB})\cdot(\vec{AC} - \vec{AB}) = AB^2 - \vec{AB}\cdot\vec{AC} = 36 - (-30) = 66$ (\textbf{S}).

\textbf{\color{brandDark}\textsf{PHẦN III. TRẢ LỜI NGẮN}}

\textbf{Câu 1.} $|\vec{AB} - \vec{AD}| = |\vec{DB}| = DB$. Vì $\widehat{BAD} = 60^\circ$ nên $\triangle ABD$ đều cạnh $12 \Rightarrow DB = \mathbf{12}$.

\textbf{Câu 2.} $\vec{MN} = \vec{AN} - \vec{AM} = -\dfrac{1}{2}\vec{AB} + \dfrac{3}{4}\vec{AC}$. Vậy $T = 8\left(\dfrac{3}{4}\right) - 2\left(-\dfrac{1}{2}\right) = 6 + 1 = \mathbf{7}$.

\textbf{Câu 3.} $(1 + 3 - 2)M = A + 3B - 2C \Rightarrow 2M = (8; 4) \Rightarrow M(4; 2) \Rightarrow T = 2(4) + 2 = \mathbf{10}$.

\textbf{Câu 4.} $\cos(\vec{u}, \vec{v}) = \dfrac{4\cdot 8 + (-3)\cdot 6}{5\cdot 10} = \dfrac{14}{50} = \mathbf{0{,}28}$.

\textbf{Câu 5.} $MA^2 = MB^2 \Leftrightarrow (m + 1)^2 + 4 = (m - 5)^2 + 16 \Leftrightarrow 12m = 36 \Leftrightarrow m = \mathbf{3}$.

\textbf{Câu 6.} $|\vec{F}| = \sqrt{300^2 + 500^2 + 2\cdot 300\cdot 500\cos 60^\circ} = \sqrt{490\,000} = \mathbf{700}\text{ N}$.

\textbf{\color{brandDark}\textsf{PHẦN IV. TỰ LUẬN}}

\textbf{Bài 1.}
a) $\vec{AM} = \dfrac{1}{2}\vec{AI} = \dfrac{1}{4}\vec{AB} + \dfrac{1}{4}\vec{AC} \Rightarrow \vec{BM} = \vec{AM} - \vec{AB} = -\dfrac{3}{4}\vec{AB} + \dfrac{1}{4}\vec{AC}$.\\
Vì $NC = 2NA \Rightarrow \vec{AN} = \dfrac{1}{3}\vec{AC} \Rightarrow \vec{BN} = \vec{AN} - \vec{AB} = -\vec{AB} + \dfrac{1}{3}\vec{AC}$.\\
b) Ta có $\vec{BM} = \dfrac{3}{4}\left(-\vec{AB} + \dfrac{1}{3}\vec{AC}\right) = \dfrac{3}{4}\vec{BN}$. Hai vectơ $\vec{BM}, \vec{BN}$ cùng phương nên $B, M, N$ thẳng hàng.

\textbf{Bài 2.}
a) $\vec{AB} = (-3; -6)$, $\vec{AC} = (2; -1) \Rightarrow \vec{AB}\cdot\vec{AC} = -6 + 6 = 0 \Rightarrow \triangle ABC$ vuông tại $A$.\\
$ABCD$ là hình bình hành $\Leftrightarrow \vec{AD} = \vec{BC} = (5; 5) \Rightarrow D = (2 + 5; 5 + 5) = (7; 10)$.\\
b) Gọi $K$ là điểm thỏa mãn $\vec{KA} + \vec{KB} + 2\vec{KC} = \vec{0}$:\\
$K = \dfrac{A + B + 2C}{4} = \left(\dfrac{2 - 1 + 8}{4}; \dfrac{5 - 1 + 8}{4}\right) = \left(\dfrac{9}{4}; 3\right)$.\\
Khi đó $T = |\vec{PA} + \vec{PB} + 2\vec{PC}| = |4\vec{PK}| = 4PK$. Vì $P \in Ox$, $PK$ nhỏ nhất khi $P$ là hình chiếu vuông góc của $K$ lên $Ox$, tức $P\!\left(\dfrac{9}{4}; 0\right)$. GTNN bằng $4 \cdot 3 = 12$.

\textbf{Bài 3.}
a) Vận tốc dòng nước: $\vec{v}_n = (3; 0)$. Đặt $\vec{v}_p = (x; y)$ với $x^2 + y^2 = 5^2 = 25$ ($y > 0$).\\
Ta có $\vec{v} = \vec{v}_p + \vec{v}_n = (x + 3; y)$. Để $\vec{v}$ cùng hướng với tia $Oy$ thì $x + 3 = 0 \Rightarrow x = -3$, suy ra $y = \sqrt{25 - 9} = 4$. Vậy $\vec{v}_p = (-3; 4)\text{ m/s}$ và $\vec{v} = (0; 4)\text{ m/s}$.\\
b) Thời gian phà sang sông: $t = \dfrac{240}{4} = 60\text{ s}$ ($1$ phút).\\
Góc $\alpha$ hợp bởi $\vec{v}_p = (-3; 4)$ và tia $Ox$ ($\vec{i} = (1; 0)$): $\cos\alpha = \dfrac{-3}{5} = -0{,}6 \Rightarrow \alpha \approx 126{,}9^\circ$.

\end{multicols}
\end{document}
